%%%%%%%%%%%%%%%%%%%%%%%%%%%%%%%%%%%%%%%%%%%%%%%%%%%%%%%%%%%%%%%%%%
%% Toto je inkludovaný LaTeXový soubor `cp_pream.tex'.          %%
%% ------------------------------------------------------------ %%
%% Verze 2022-06-03                                             %%
%% Vyvíjí & udržuje <stanislav.hledik@physics.slu.cz>           %%
%%%%%%%%%%%%%%%%%%%%%%%%%%%%%%%%%%%%%%%%%%%%%%%%%%%%%%%%%%%%%%%%%%

\chapter{Preambule}\label{preamble}%%%%%%%%%%%%%%%%%%%%%%%%%%%%%%%%%%%%%%%%

Úvodní část \LaTeX ového hlavního zdrojového souboru nacházející se mezi
specifikací dokumentní třídy v~\pgm{\textbackslash
  documentclass[}$\langle$\emph{volby}$\rangle$\pgm{]\{\ipotname\}} a začátkem
prostředí \verb!\begin{document}! se nazývá
  \emph{preambule}~\cite{Ryb:2005:LaTeX:3}. Umisťují se do ní příkazy pro
  načítání balíčků pomocí příkazu \pgm{\textbackslash usepackage} a
  uživatelská nastavení, definice či redefinice. Po \verb!\begin{document}!
    následuje \emph{tělo dokumentu}, kam se zapisuje zejména vlastní text
s~potřebným kontextovým značkováním.

U~dokumentní třídy \pgm{\ipotname .cls} je doporučeno umisťovat do preambule
pouze načítání balíčků, avšak uživatelská nastavení, definice a redefinice
raději přesunout do konfiguračního souboru \pgm{\ipotname .cfg} popsaného
v~Dodatku~\ref{config}. Primárním důvodem tohoto doporučení je, že
konfigurační soubor \pgm{\ipotname .cfg} je načítán v~argumentu makra
\pgm{\textbackslash AtEndPreamble}~\cite{Wri-Leh:LIVE:CTAN:etoolbox} až po
balíčku \pkg{hyperref}, jenž sám redefinuje některé příkazy (viz
Dodatek~\ref{packages}). Sekundární důvod je víceméně estetický~-- sleduje
oddělení souboru s~konfigurací a souboru s~textem. Jinak ale lze vaše definice
a redefinice v~preambuli uvést, pokud si je otestujete.

Celý hlavní zdrojový soubor \pgm{\thesnamecze .tex} tohoto manuálu jakožto
česky psané fiktivní závěrečné práce je včetně preambule bohatě komentován.
Doporučujeme čtenáři tyto komentáře projít; ostatně vezme-li soubor
\pgm{\thesnamecze .tex} jako základ pro psaní vlastní práce, je četba
komentářů zcela dostačující. Níže uvedené fragmenty preambule jsou načteny
přímo z~\pgm{\thesnamecze .tex} pomocí balíčku
\pkg{fancyvrb}~\cite{Vos:LIVE:CTAN:fancyvrb}.

\section{Specifikace dokumentní třídy}\label{docclass}%%%%%%%%%%%%%%%%%%%%%

Preambuli samotné předchází zavedení dokumentní třídy \pgm{\ipotname .cls} a
nastavení jejích voleb, jež jsou popsány v~Dodatku~\ref{options}. Právě jedna
z~voleb \pgm{draft} nebo \pgm{final} musí zůstat odkomentovaná (nenechávajte
je obě zakomentované nebo odkomentované); ostatní voleb lze použít podle
vašich vlastních požadavků (implicitní chování při zakomentovaní volbě je
uvedeno v~závorce komentáře):

\VerbatimInput[frame=none,firstline=16,lastline=24,fontsize=\small]%
{vzorova_prace.tex}

\section{Fonty a jazyková lokalizace}\label{fontslang}%%%%%%%%%%%%%%%%%%%%%

Nyní začíná vlastní preambule. Pro zajištění vyšší flexibility nejsou
příslušné balíčky načítány v~dokumentní třídě (viz Dodatek~\ref{packages}),
ale až zde v~preambuli.

Balíčky načtené v~následujícím kódu způsobí globální použití nadrodiny fontů
Latin Modern (\pkg{lm}~\cite{Jac:LIVE:CTAN:lm}) místo nadrodiny fontů Computer
Modern (\pkg{cm}). Je nutné uvést kódování fontu volbou \opt{T1} balíčku
\pkg{fontenc}~\cite{LtX:LIVE:CTAN:fontenc} a kódování zdrojových textů volbou
\opt{utf8} balíčku \pkg{inputenc}~\cite{LtX:LIVE:CTAN:inputenc}, jež je dnes
standardem na všech dnešních hlavních operačních systémech:\footnote{Nebylo
  tomu tak vždy. Případná jiná kódování, která připadají v~úvahu, jsou
  \opt{cp1250} na Windows nebo \opt{latin2} na Linuxu, ale byla by nutná
  konverze.}

\VerbatimInput[frame=none,firstline=31,lastline=33,fontsize=\small]%
{vzorova_prace.tex}

Následuje blok jazykové lokalizace. Závěrečná práce v~češtině \emph{musí} mít
nastavenu češtinu jako hlavní jazyk volbou \opt{czech} uvedenou na
\emph{posledním místě} v~seznamu volitelných argumentů balíčku
\pkg{babel}. Jako další jazyk \emph{musí} být volbou \opt{english} zavedena
angličtina. Kromě těchto dvou povinných jazyků lze zavést i~další jazyky;
aktuálně jsou v~\pkg{\ipotname} podporovány: \opt{slovak}, \opt{french},
\opt{german} (viz též dokumentaci balíčku
\pkg{babel}~\cite{Bra:LIVE:CTAN:Babel}):

\VerbatimInput[frame=none,firstline=42,lastline=42,fontsize=\small]%
{vzorova_prace.tex}

Preferujete-li Times jako textový i~matematický font\footnote{Font Times je
  prostorově úspornější. Nepoužívejte balíček
  \pkg{times}~\cite{Rah:LIVE:CTAN:times}, který změní jen textový, nikoli však
  matematický font, ani balíčky \pkg{mathptmx} nebo
  \pkg{psnfss}~\cite{Vie:LIVE:CTAN:mathptmx,%
    Har:LIVE:CTAN:psnfss}.}, stačí odkomentovat následující řádek pro načtení
balíčku \pkg{txfonts}~\cite{Ryu:LIVE:CTAN:txfonts}. Necháte-li jej
zakomentovaný, bude použit Latin Modern, jehož zavedení je popsáno výše.

\VerbatimInput[frame=none,firstline=46,lastline=46,fontsize=\small]%
{vzorova_prace.tex}

\noindent
Použijete-li výše uvedený balíček \pkg{txfonts}~\cite{Ryu:LIVE:CTAN:txfonts},
bezpatkový (sans serif) font a strojopisný (monospace) font budou změněny po
řadě na fonty Helvetica a Courier. Pro zachování bezpatkového a strojopisného
fontu z~nadrodiny \pkg{lm}~\cite{Jac:LIVE:CTAN:lm}, což podle názoru autora
dává lepší estetický výsledek, stačí odkomentovat tyto čtyři řádky (nechejte
je zakomentované, je-li zakomentováno přecházející načtení balíčku
\pkg{txfonts}):

\VerbatimInput[frame=none,firstline=53,lastline=56,fontsize=\small]%
{vzorova_prace.tex}

\section{Podpora tvorby rejstříku}\label{indexsupp}%%%%%%%%%%%%%%%%%%%%%%%%

Následující kód odkomentují a balíček
\pkg{makeidx}~\cite{LtX:LIVE:CTAN:makeidx,%
  Bee:LIVE:CTAN:MakeIndex} nechají načíst pouze ti, kdo chtějí vytvářet
rejstřík. Jelikož v~tomto manuálu rejstřík není, odkazujeme na další informace
např. v~monografiích~\cite{Ryb:2005:LaTeX:3,%
  Kop-Dal:2004:LaTeXPodrPruv:} a dokumentaci~\cite{LtX:LIVE:CTAN:makeidx,%
  Bee:LIVE:CTAN:MakeIndex}:

\VerbatimInput[frame=none,firstline=59,lastline=59,fontsize=\small]%
{vzorova_prace.tex}

\section{Další balíčky}\label{otherpack}%%%%%%%%%%%%%%%%%%%%%%%%%%%%%%%%%%%

Dále můžete podle potřeby načíst další balíčky kompatibilní s~dokumentní
třídou \pkg{\ipotname}. Doporučuje se načítat jen balíčky, které jsou opravdu
nezbytné. Balíčky uvedené v~Dodatku~\ref{packages} a jejich závislosti se
načítají v~dokumentní třídě a je tedy zbytečné a škodlivé je zde načítat. Pro
tento manuál potřebujeme balíček \pkg{fancyvrb}~\cite{Vos:LIVE:CTAN:fancyvrb}
umožňující načítání fragmentů kódu přímo ze zdrojových souborů jako v~této
kapitole; další balíček \pkg{subfig}~\cite{biditex:LIVE:CTAN:subfig} je
připraven k~použití:

\VerbatimInput[frame=none,firstline=66,lastline=67,fontsize=\small]%
{vzorova_prace.tex}

\section{Nastavení, definice a redefinice}\label{settdefredef}%%%%%%%%%%%%%

Odsud až po \verb!\begin{document}! lze umístit vaše privátní nastavení,
  definice a redefinice. Je však výhodnější umístit je do konfiguračního
  souboru \pgm{\ipotname .cfg} popsaného v~Dodatku~\ref{config}, jenž bude
  automaticky (je-li přítomen) nalezen a načten~-- důvody jsou vysvětleny
  v~úvodu této kapitoly. Některá nastavení však budou fungovat jen když jsou
  v~konfiguračním souboru, např. nastavení \pgm{\textbackslash
    pdfcompresslevel}.  Z~toho důvodu se doporučuje upřednostnit konfigurační
  soubor. Některá nastavení budou fungovat i~odtud (aktivní jsou však
  v~\pgm{\ipotname .cfg}), jako například:

\VerbatimInput[frame=none,firstline=79,lastline=81,fontsize=\small]%
{vzorova_prace.tex}

\section{Vzory dělení slov}\label{userhyph}%%%%%%%%%%%%%%%%%%%%%%%%%%%%%%%%

Uživatelské vzory dělení slov musí být umístěny zde v~preambuli, nikoli
v~konfiguračním souboru \pgm{\ipotname .cfg}:

\VerbatimInput[frame=none,firstline=85,lastline=91,fontsize=\small]%
{vzorova_prace.tex}

\section{Konec preambule}\label{preambend}%%%%%%%%%%%%%%%%%%%%%%%%%%%%%%%%%

Začátkem prostředí \pgm{document} končí preambule:

\VerbatimInput[frame=none,firstline=97,lastline=97,fontsize=\small]%
{vzorova_prace.tex}

\endinput

%% 
%% End of file `cp_pream.tex'.
